\documentclass{beamer}
\usetheme{Madrid}
\usecolortheme{default}
\usepackage{amsmath, amssymb, graphicx, hyperref, array, booktabs, multicol}
\usepackage[utf8]{inputenc}
\usepackage{color}

% Define colors
\definecolor{cfablue}{RGB}{0,51,102}
\definecolor{cfagold}{RGB}{204,153,0}
\definecolor{cfared}{RGB}{180,0,0}
\definecolor{cfagreen}{RGB}{0,120,0}
\setbeamercolor{title}{fg=cfablue}
\setbeamercolor{frametitle}{fg=cfablue}
\setbeamercolor{structure}{fg=cfablue}

% REDUCED FONT SIZE FOR ALL SLIDES
\renewcommand{\normalsize}{\fontsize{8pt}{10pt}\selectfont}
\renewcommand{\small}{\fontsize{7pt}{9pt}\selectfont}
\renewcommand{\footnotesize}{\fontsize{6pt}{8pt}\selectfont}
\renewcommand{\tiny}{\fontsize{5pt}{7pt}\selectfont}


% Title information
\title{The Relationship Between Legality and Ethics in the Finance Profession}
\subtitle{CFA Institute Standards of Practice Handbook Analysis}
\author{Scholar Practitioner Candidate}
\date{\today}

\begin{document}
	
	% ============================================================
	% SLIDE 1: TITLE
	% ============================================================
	\begin{frame}
		\titlepage
	\end{frame}
	
	% ============================================================
	% SLIDE 2: AGENDA
	% ============================================================
	\begin{frame}{Agenda}
		\begin{enumerate}
			\item \textbf{Core Question}: Does Illegal Always Mean Unethical?
			\item \textbf{CFA Framework}: The "More Strict" Standard
			\item \textbf{Young's Argument}: Valid Points or Justification?
			\item \textbf{Challenging Scenarios} (20 Very Hard Questions)
			\item \textbf{Key Takeaways} for DBA Scholars
		\end{enumerate}
	\end{frame}
	
	% ============================================================
	% SLIDE 3: THE CENTRAL TENSION
	% ============================================================
	\begin{frame}{The Central Tension}
		\begin{block}{The Question}
			Does illegal always mean unethical? More importantly, does legal always mean ethical?
		\end{block}
		
		\begin{columns}
			\column{0.5\textwidth}
			\textbf{Common Assumption}
			\begin{itemize}
				\item Illegal = Unethical
				\item Legal = Ethical
				\item Compliance = Sufficient
			\end{itemize}
			
			\column{0.5\textwidth}
			\textbf{CFA Position}
			\begin{itemize}
				\item Not necessarily
				\item Legal is baseline, not ceiling
				\item \textcolor{cfablue}{"More strict" standard applies}
			\end{itemize}
		\end{columns}
		
		\vspace{0.3cm}
		\footnotesize \textit{Source: CFA Institute Standards of Practice Handbook (2014, p. 14)}
	\end{frame}
	
	% ============================================================
	% SLIDE 4: STANDARD I(A) - KNOWLEDGE OF THE LAW
	% ============================================================
	\begin{frame}{Standard I(A): Knowledge of the Law}
		\begin{block}{The Standard}
			"Members and Candidates must understand and comply with all applicable laws, rules, and regulations... In the event of conflict, Members and Candidates must comply with the more strict law, rule, or regulation."
		\end{block}
		
		\begin{itemize}
			\item \textbf{Key Principle}: Legal compliance is the \textcolor{cfablue}{minimum}
			\item \textbf{Global Application}: Different jurisdictions, different rules
			\item \textbf{The "More Strict" Test}: Compare law vs. Code \& Standards
		\end{itemize}
		
		\vspace{0.3cm}
		\begin{center}
			\fcolorbox{cfablue}{white}{\parbox{0.8\textwidth}{\centering \textbf{When in conflict, follow the HIGHER standard}}}
		\end{center}
	\end{frame}
	
	% ============================================================
	% SLIDE 5: GLOBAL APPLICATION - EXHIBIT 1
	% ============================================================
	\begin{frame}{Global Application of the Code and Standards}
		\begin{table}[]
			\centering
			\small
			\begin{tabular}{|p{3.2cm}|p{3.5cm}|}
				\hline
				\textbf{Applicable Law} & \textbf{Duty} \\
				\hline
				No securities laws & Adhere to Code \& Standards \\
				Less strict than CFA & Adhere to Code \& Standards \\
				More strict than CFA & Adhere to local law \\
				\hline
			\end{tabular}
		\end{table}
		
		\vspace{0.3cm}
		\textbf{Implication}: Legal requirements vary; ethical obligations are universal.
		
		\vspace{0.2cm}
		\footnotesize \textit{Source: CFA Institute (2014, pp. 16-18)}
	\end{frame}
	
	% ============================================================
	% SLIDE 6: YOUNG'S ARGUMENT - OVERVIEW
	% ============================================================
	\begin{frame}{Young's Argument (2019)}
		\begin{block}{Core Thesis}
			"Putting the law in its place" - Challenging the assumption that illegal implies unethical.
		\end{block}
		
		\textbf{Key Points:}
		\begin{enumerate}
			\item Laws can be \textcolor{cfablue}{unjust} or \textcolor{cfablue}{outdated}
			\item Regulatory bodies can be flawed ("patronage haven")
			\item Civil disobedience is ethically defensible
			\item Legal is not Moral; Illegal is not Immoral
		\end{enumerate}
		
		\vspace{0.3cm}
		\textbf{Critical Question}: Is Young exposing flawed ethical reasoning or justifying misconduct?
	\end{frame}
	
	% ============================================================
	% SLIDE 7: HARD QUESTION 1 - MULTI-JURISDICTION CONFLICT
	% ============================================================
	\begin{frame}{Hard Question 1: Multi-Jurisdiction Compliance Dilemma}
		\begin{block}{The Scenario (CFA Handbook, pp. 16-18, Exhibit 1)}
			Laura Jameson works for a US-based multinational investment adviser and is registered in Karramba, a wealthy island nation. Karramban law PROHIBITS registered advisers from participating in IPOs for personal accounts. US law has NO such prohibition. Jameson participates in an IPO for her personal account, believing she should follow US law as a US citizen.
		\end{block}
		
		\vspace{0.2cm}
		\textbf{Questions:}
		\begin{enumerate}
			\item Did Jameson violate Standard I(A)? \textcolor{cfared}{[Answer: YES]}
			\item Which law applies - US or Karramban? \textcolor{cfared}{[Answer: Karramban]}
			\item If Jameson is a CFA charterholder, does that change the analysis? \textcolor{cfared}{[Answer: No]}
			\item What if the local law is LESS strict than the Code and Standards? \textcolor{cfared}{[Answer: Follow Code \& Standards]}
		\end{enumerate}
		
		\vspace{0.2cm}
		\textcolor{cfablue}{\textbf{CFA Answer}}: Violates Standard I(A). As a registered adviser in Karramba, she must follow local law. The "more strict" standard applies - if local law is stricter, follow local law.
		
		\vspace{0.1cm}
		\textcolor{cfared}{\textbf{Common Mistake}}: Thinking US citizenship or CFA designation exempts one from local laws. It does NOT.
	\end{frame}
	
	% ============================================================
	% SLIDE 8: HARD QUESTION 2 - UNETHICAL BUT LEGAL CONDUCT
	% ============================================================
	\begin{frame}{Hard Question 2: Legal But Unethical Conduct}
		\begin{block}{The Scenario (CFA Handbook, p. 55, Standard I(D))}
			Simon Sasserman, a trust investment officer, drinks heavily at lunch daily. His clients observe him at the country club consuming multiple drinks. Back at work, he is intoxicated while making investment decisions. His colleagues handle business with him only in the morning because they distrust his judgment after lunch. His conduct is NOT illegal.
		\end{block}
		
		\vspace{0.2cm}
		\textbf{Questions:}
		\begin{enumerate}
			\item Does Sasserman violate any CFA Standard? \textcolor{cfared}{[Answer: YES - Standard I(D)]}
			\item If conduct is legal, can it still be a violation? \textcolor{cfared}{[Answer: YES]}
			\item What if he drinks only after hours? \textcolor{cfared}{[Answer: May not violate if it doesn't affect professional conduct]}
			\item Would personal bankruptcy violate this standard? \textcolor{cfared}{[Answer: Only if involves fraud/deceit]}
		\end{enumerate}
		
		\vspace{0.2cm}
		\textcolor{cfablue}{\textbf{CFA Answer}}: Violates Standard I(D) - Misconduct. Conduct reflects adversely on professional reputation, integrity, and competence. The standard covers conduct that damages trustworthiness or competence, even if not illegal.
		
		\vspace{0.1cm}
		\textcolor{cfared}{\textbf{Common Mistake}}: Assuming only illegal conduct violates the Standards. Illegal behavior is NOT required for a violation.
	\end{frame}
	
	% ============================================================
	% SLIDE 9: HARD QUESTION 3 - DISSOCIATION AND WHISTLEBLOWING
	% ============================================================
	\begin{frame}{Hard Question 3: Dissociation from Unethical Activity}
		\begin{block}{The Scenario (CFA Handbook, pp. 20-21, Example 3)}
			Kamisha Washington discovers her firm's advertised 10-year performance composite OMITS accounts that left the firm, inflating returns. The promotional material contains this erroneous performance number. She is asked to use this material to solicit new clients.
		\end{block}
		
		\vspace{0.2cm}
		\textbf{Questions:}
		\begin{enumerate}
			\item What must Washington do? \textcolor{cfared}{[Answer: Dissociate from activity]}
			\item If she DID NOT calculate the performance herself, is she responsible? \textcolor{cfared}{[Answer: YES]}
			\item What intermediate steps should she take? \textcolor{cfared}{[Answer: Bring to supervisor/compliance]}
			\item If firm insists she use the material, what should she consider? \textcolor{cfared}{[Answer: Seek other employment]}
		\end{enumerate}
		
		\vspace{0.2cm}
		\textcolor{cfablue}{\textbf{CFA Answer}}: Washington would be assisting in violation if she uses inflated numbers. Must dissociate - bring to supervisor/compliance, refuse to use material, consider leaving if firm insists.
		
		\vspace{0.1cm}
		\textcolor{cfared}{\textbf{Common Mistake}}: Thinking "I didn't create the misrepresentation so I'm not responsible." Knowledge and participation matter.
	\end{frame}
	
	% ============================================================
	% SLIDE 10: HARD QUESTION 4 - SOCIAL MEDIA COMPLIANCE
	% ============================================================
	\begin{frame}{Hard Question 4: Social Media and Regulatory Knowledge}
		\begin{block}{The Scenario (CFA Handbook, p. 23, Example 8)}
			Colleen White posts investment information, performance reports, and investment opinions to her Facebook page. She sends tweets like "Prospects for future growth of XYZ company look good! \#makingmoney4U." Prior to this, the local regulator issued new requirements for online electronic communication. Her communications conflict with the recent regulatory announcements.
		\end{block}
		
		\vspace{0.2cm}
		\textbf{Questions:}
		\begin{enumerate}
			\item Did White violate Standard I(A)? \textcolor{cfared}{[Answer: YES]}
			\item Does social media change the compliance requirements? \textcolor{cfared}{[Answer: No]}
			\item What if she didn't know about the new regulations? \textcolor{cfared}{[Answer: Still a violation - must stay informed]}
			\item What sources should she consult? \textcolor{cfared}{[Answer: Compliance, external counsel, service providers]}
		\end{enumerate}
		
		\vspace{0.2cm}
		\textcolor{cfablue}{\textbf{CFA Answer}}: Violates Standard I(A) because communications don't comply with existing guidance. Must be aware of evolving legal requirements for new areas.
		
		\vspace{0.1cm}
		\textcolor{cfared}{\textbf{Common Mistake}}: Assuming social media is exempt from regulations or that "I didn't know" is a defense.
	\end{frame}
	
	% ============================================================
	% SLIDE 11: HARD QUESTION 5 - DIRECTOR CONFLICT
	% ============================================================
	\begin{frame}{Hard Question 5: Director Conflict of Interest}
		\begin{block}{The Scenario (CFA Handbook, pp. 191-192, Example 8)}
			Carol Corky, senior portfolio manager at Universal Management (US\$100M AUM), becomes a trustee of the Chelsea Foundation (US\$2B AUM). Chelsea is a large not-for-profit foundation. Corky believes informing Universal is not necessary. Universal manages only individual investor accounts.
		\end{block}
		
		\vspace{0.2cm}
		\textbf{Questions:}
		\begin{enumerate}
			\item Did Corky violate Standard VI(A)? \textcolor{cfared}{[Answer: YES]}
			\item What conflicts does board service create? \textcolor{cfared}{[Answer: Time conflict, material nonpublic info, client duties vs. shareholder duties]}
			\item If Corky does NOT manage Chelsea's investments, is disclosure still required? \textcolor{cfared}{[Answer: YES]}
			\item What if Universal has no policy on outside board service? \textcolor{cfared}{[Answer: Still must disclose]}
		\end{enumerate}
		
		\vspace{0.2cm}
		\textcolor{cfablue}{\textbf{CFA Answer}}: Violates Standard VI(A). Must disclose because board service can be time-consuming and may create conflicts with portfolio responsibilities. May also become involved in Chelsea's investment decisions.
		
		\vspace{0.1cm}
		\textcolor{cfared}{\textbf{Common Mistake}}: Thinking only monetary compensation matters. Time commitment and potential access to nonpublic information create conflicts.
	\end{frame}
	
	% ============================================================
	% SLIDE 12: HARD QUESTION 6 - MATERIAL NONPUBLIC INFO COMPLEXITY
	% ============================================================
	\begin{frame}{Hard Question 6: Expert Network and Material Nonpublic Info}
		\begin{block}{The Scenario (CFA Handbook, p. 73, Example 12)}
			Tom Watson, hedge fund research analyst, arranges a call with a scientist leading tests on a new semiconductor product. The scientist indicates disappointment with performance. Watson relays insights to his fund. The fund sells its position and buys put options because the market expects success. The company's share price reflects market optimism.
		\end{block}
		
		\vspace{0.2cm}
		\textbf{Questions:}
		\begin{enumerate}
			\item Did Watson violate Standard II(A)? \textcolor{cfared}{[Answer: YES]}
			\item What if the scientist signed an agreement not to disclose confidential information? \textcolor{cfared}{[Answer: Watson still responsible for his own determination]}
			\item Does the expert network's compliance procedures protect Watson? \textcolor{cfared}{[Answer: No - must make own determination]}
			\item What is the materiality threshold? \textcolor{cfared}{[Answer: Ongoing product test results affecting market expectations]}
		\end{enumerate}
		
		\vspace{0.2cm}
		\textcolor{cfablue}{\textbf{CFA Answer}}: Violates Standard II(A). Watson cannot rely on agreements signed by experts. Must make own determination whether information is material nonpublic.
		
		\vspace{0.1cm}
		\textcolor{cfared}{\textbf{Common Mistake}}: Assuming expert networks or third-party agreements protect against insider trading violations. They do NOT.
	\end{frame}
	
	% ============================================================
	% SLIDE 13: HARD QUESTION 7 - MODEL INPUT MANIPULATION
	% ============================================================
	\begin{frame}{Hard Question 7: Model Input Manipulation}
		\begin{block}{The Scenario (CFA Handbook, pp. 79-80, Example 8)}
			Bill Mandeville supervises structured financing. To achieve the best rating, he uses mostly positive scenarios - reflecting minimal downside risk in underlying assets. Output statistics support minimal potential downside risk. Compensation is partially based on rating level and successful sale of products, NOT linked to long-term performance. The economy later declines, causing defaults and major turmoil.
		\end{block}
		
		\vspace{0.2cm}
		\textbf{Questions:}
		\begin{enumerate}
			\item Did Mandeville violate Standard II(B)? \textcolor{cfared}{[Answer: YES]}
			\item What if the model outputs were technically accurate given the inputs? \textcolor{cfared}{[Answer: Intent matters - selecting inputs to achieve desired rating is manipulation]}
			\item Does the fact that the product was legal make it ethical? \textcolor{cfared}{[Answer: No]}
			\item What broader damage did Mandeville's actions cause? \textcolor{cfared}{[Answer: Loss of trust in structured products and capital markets]}
		\end{enumerate}
		
		\vspace{0.2cm}
		\textcolor{cfablue}{\textbf{CFA Answer}}: Violates Standard II(B). Manipulating inputs to minimize risk misleads market participants, damages trust in capital markets, and demonstrates intent to mislead.
		
		\vspace{0.1cm}
		\textcolor{cfared}{\textbf{Common Mistake}}: Thinking legal = ethical or that "technically accurate" inputs justify manipulation. Intent is critical.
	\end{frame}
	
	% ============================================================
	% SLIDE 14: HARD QUESTION 8 - CIVIL DISOBEDIENCE
	% ============================================================
	\begin{frame}{Hard Question 8: Civil Disobedience and Professional Conduct}
		\begin{block}{The Scenario (CFA Handbook, p. 57, Example 4)}
			Carmen Garcia manages a socially responsible investing mutual fund. She is also an environmental activist. As a result of participating in nonviolent protests, she has been arrested numerous times for trespassing on the property of a petrochemical plant accused of damaging the environment.
		\end{block}
		
		\vspace{0.2cm}
		\textbf{Questions:}
		\begin{enumerate}
			\item Does Garcia violate Standard I(D)? \textcolor{cfared}{[Answer: Generally NO]}
			\item Why would civil disobedience NOT violate the standard? \textcolor{cfared}{[Answer: Doesn't reflect poorly on professional reputation/integrity/competence]}
			\item What if the protests were violent? \textcolor{cfared}{[Answer: Could violate if reflects poorly on professional competence]}
			\item What if clients become uncomfortable with her activism? \textcolor{cfared}{[Answer: Firm policies may address, but not a Standards violation]}
		\end{enumerate}
		
		\vspace{0.2cm}
		\textcolor{cfablue}{\textbf{CFA Answer}}: Generally, Standard I(D) does not cover legal transgressions from civil disobedience in support of personal beliefs. Such conduct does not reflect poorly on professional reputation, integrity, or competence.
		
		\vspace{0.1cm}
		\textcolor{cfared}{\textbf{Common Mistake}}: Assuming ANY arrest violates Standard I(D). Context and professional relevance matter.
	\end{frame}
	
	% ============================================================
	% SLIDE 15: HARD QUESTION 9 - COMPENSATION STRUCTURES
	% ============================================================
	\begin{frame}{Hard Question 9: Compensation and Client Interest Conflicts}
		\begin{block}{The Scenario (CFA Handbook, p. 190, Example 5)}
			Samantha Snead manages long-term public retirement funds and pension plans. Her employer introduces a bonus compensation system rewarding portfolio managers based on QUARTERLY performance relative to peers and benchmarks. Snead changes strategy, purchasing high-beta stocks seemingly contrary to clients' IPSs. A client asks why the portfolio is dominated by high-beta stocks. No change in objective or strategy was recommended.
		\end{block}
		
		\vspace{0.2cm}
		\textbf{Questions:}
		\begin{enumerate}
			\item Did Snead violate Standard VI(A)? \textcolor{cfared}{[Answer: YES]}
			\item Did she violate Standard III(C)? \textcolor{cfared}{[Answer: YES]}
			\item If the bonus was ANNUAL instead of QUARTERLY, would that change the analysis? \textcolor{cfared}{[Answer: Less immediate pressure, but still must disclose conflicts]}
			\item What is Snead required to disclose? \textcolor{cfared}{[Answer: Compensation arrangement changes that create conflicts]}
		\end{enumerate}
		
		\vspace{0.2cm}
		\textcolor{cfablue}{\textbf{CFA Answer}}: Violates Standard VI(A) by failing to inform clients of changes in compensation creating conflicts with IPSs. Also violates Standard III(C) - suitability. Short-term performance goals conflict with long-term account objectives.
		
		\vspace{0.1cm}
		\textcolor{cfared}{\textbf{Common Mistake}}: Thinking compensation arrangements don't need disclosure if "everyone does it." Material conflicts MUST be disclosed.
	\end{frame}
	
	% ============================================================
	% SLIDE 16: HARD QUESTION 10 - GIFT ACCEPTANCE LIMITS
	% ============================================================
	\begin{frame}{Hard Question 10: Gift Acceptance and Independence}
		\begin{block}{The Scenario (CFA Handbook, p. 35, Example 7)}
			Edward Grant directs large commission business to a New York brokerage house. In appreciation, the brokerage gives Grant two tickets to the World Cup in South Africa, two nights at a resort, meals, and limousine transportation. Grant fails to disclose this package to his supervisor.
		\end{block}
		
		\vspace{0.2cm}
		\textbf{Questions:}
		\begin{enumerate}
			\item Did Grant violate Standard I(B)? \textcolor{cfared}{[Answer: YES]}
			\item What if the value of the gifts was modest? \textcolor{cfared}{[Answer: Must still consider whether it could compromise objectivity]}
			\item What if Grant disclosed the gifts to his supervisor? \textcolor{cfared}{[Answer: Might still be violation if the gifts could compromise independence]}
			\item What is the best practice regarding gifts? \textcolor{cfared}{[Answer: Reject anything that could threaten independence; token items only]}
		\end{enumerate}
		
		\vspace{0.2cm}
		\textcolor{cfablue}{\textbf{CFA Answer}}: Violates Standard I(B). Accepting substantial gifts may impede independence and objectivity. Must avoid situations that could cause or be perceived to cause loss of objectivity.
		
		\vspace{0.1cm}
		\textcolor{cfared}{\textbf{Common Mistake}}: Thinking disclosure makes ANY gift acceptable. Gifts that could compromise objectivity are prohibited regardless of disclosure.
	\end{frame}
	
	% ============================================================
	% SLIDE 17: HARD QUESTION 11 - PAY-TO-PLAY SCANDAL
	% ============================================================
	\begin{frame}{Hard Question 11: Pay-to-Play and Manager Selection}
		\begin{block}{The Scenario (CFA Handbook, p. 39, Example 13)}
			Adrian Mandel, CFA, makes significant monetary contributions to Ernesto Gomez's union reelection campaign. Gomez chairs the investment board of UDPBU pension fund. Mandel also entertains Gomez lavishly. Expenses are reimbursed as "marketing expenses" and disclosed to senior management.
		\end{block}
		
		\vspace{0.2cm}
		\textbf{Questions:}
		\begin{enumerate}
			\item Did Mandel violate Standard I(B)? \textcolor{cfared}{[Answer: YES]}
			\item If political contributions are legal, can they still be unethical? \textcolor{cfared}{[Answer: YES]}
			\item Does disclosure to senior management make it acceptable? \textcolor{cfared}{[Answer: No - still violates independence]}
			\item How does this differ from legitimate marketing? \textcolor{cfared}{[Answer: Intent to compromise objectivity]}
		\end{enumerate}
		
		\vspace{0.2cm}
		\textcolor{cfablue}{\textbf{CFA Answer}}: Violates Standard I(B). Mandel gave gifts and benefits that could reasonably compromise Gomez's objectivity. The conduct is unethical even if structured to be technically legal.
		
		\vspace{0.1cm}
		\textcolor{cfared}{\textbf{Common Mistake}}: Assuming legal structuring or internal disclosure eliminates ethical violations. Intent and impact on objectivity matter.
	\end{frame}
	
	% ============================================================
	% SLIDE 18: HARD QUESTION 12 - FAIR DEALING WITH REFERRALS
	% ============================================================
	\begin{frame}{Hard Question 12: Fair Dealing and Referral Arrangements}
		\begin{block}{The Scenario (CFA Handbook, p. 204, Example 2)}
			James Handley works in Central Trust Bank's trust department. He receives compensation for each referral to the brokerage and personal financial management departments that results in a sale. He refers several clients but does not disclose the arrangement.
		\end{block}
		
		\vspace{0.2cm}
		\textbf{Questions:}
		\begin{enumerate}
			\item Did Handley violate Standard VI(C)? \textcolor{cfared}{[Answer: YES]}
			\item Does Standard VI(C) distinguish between internal and external payments? \textcolor{cfared}{[Answer: No - both must be disclosed]}
			\item When must disclosure be made? \textcolor{cfared}{[Answer: At time of referral]}
			\item What must the disclosure include? \textcolor{cfared}{[Answer: Nature and value of benefit, in writing]}
		\end{enumerate}
		
		\vspace{0.2cm}
		\textcolor{cfablue}{\textbf{CFA Answer}}: Violates Standard VI(C). Standard does not distinguish between internal and external referral payments. Must disclose nature and value of benefit in writing at time of referral.
		
		\vspace{0.1cm}
		\textcolor{cfared}{\textbf{Common Mistake}}: Thinking internal referral arrangements don't need disclosure. They do.
	\end{frame}
	
	% ============================================================
	% SLIDE 19: HARD QUESTION 13 - PERFORMANCE PRESENTATION
	% ============================================================
	\begin{frame}{Hard Question 13: Performance Presentation Misrepresentation}
		\begin{block}{The Scenario (CFA Handbook, pp. 114-115, Example 1)}
			Kyle Taylor states in a brochure: "You can expect steady 25\% annual compound growth of the value of your investments over the year." Taylor Trust's common trust fund achieved 25\% for ONE year, which mirrored market increase. The fund has never averaged that growth for more than one year. Average growth of all trust accounts over five years is 5\% per year.
		\end{block}
		
		\vspace{0.2cm}
		\textbf{Questions:}
		\begin{enumerate}
			\item Did Taylor violate Standard III(D)? \textcolor{cfared}{[Answer: YES]}
			\item Did Taylor violate Standard I(C)? \textcolor{cfared}{[Answer: YES]}
			\item What if Taylor said "past performance" instead of "can expect"? \textcolor{cfared}{[Answer: May not guarantee, but still must present complete performance]}
			\item What would a proper presentation include? \textcolor{cfared}{[Answer: Disclosure of one-year anomaly, all accounts, no guarantees]}
		\end{enumerate}
		
		\vspace{0.2cm}
		\textcolor{cfablue}{\textbf{CFA Answer}}: Violates Standard III(D) - incomplete and misleading presentation. Also violates Standard I(C) - guarantees future performance. Must disclose that 25\% occurred only one year and include all accounts.
		
		\vspace{0.1cm}
		\textcolor{cfared}{\textbf{Common Mistake}}: Thinking past performance can be presented as "expected" future performance. It cannot.
	\end{frame}
	
	% ============================================================
	% SLIDE 20: HARD QUESTION 14 - RECORD RETENTION PROPERTY RIGHTS
	% ============================================================
	\begin{frame}{Hard Question 14: Record Retention and Firm Property}
		\begin{block}{The Scenario (CFA Handbook, p. 183, Example 3)}
			Martin Blank develops an analytical model while employed at GPIM. He systematically documents all assumptions and reasoning. Due to the model's success, Blank is hired as head of research at a competitor. Blank takes copies of all supporting records to his new employer without GPIM's permission.
		\end{block}
		
		\vspace{0.2cm}
		\textbf{Questions:}
		\begin{enumerate}
			\item Did Blank violate Standard V(C)? \textcolor{cfared}{[Answer: YES]}
			\item What if Blank developed the model on his own time? \textcolor{cfared}{[Answer: Still firm property if used firm resources]}
			\item What if he only takes "ideas" not actual records? \textcolor{cfared}{[Answer: May be acceptable if not using firm property]}
			\item How can Blank use the model in the future? \textcolor{cfared}{[Answer: Re-create records at new firm from public sources]}
		\end{enumerate}
		
		\vspace{0.2cm}
		\textcolor{cfablue}{\textbf{CFA Answer}}: Violates Standard V(C). Records created as part of professional activity are firm property. Blank must re-create records at the new firm from public sources.
		
		\vspace{0.1cm}
		\textcolor{cfared}{\textbf{Common Mistake}}: Thinking work done on personal time or "ideas" are personal property. Records created using firm resources are firm property.
	\end{frame}
	
	% ============================================================
	% SLIDE 21: HARD QUESTION 15 - RESEARCH INDEPENDENCE AND PRESSURE
	% ============================================================
	\begin{frame}{Hard Question 15: Research Independence and Pressure}
		\begin{block}{The Scenario (CFA Handbook, p. 34, Example 3)}
			Walter Fritz, equity analyst, concludes Metals \& Mining stock is overpriced. He is concerned a negative report will hurt the good relationship between Metals \& Mining and Hilton Brokerage's investment banking division. A senior manager sent him a proposal for a debt offering. Fritz needs to produce a report quickly.
		\end{block}
		
		\vspace{0.2cm}
		\textbf{Questions:}
		\begin{enumerate}
			\item What must Fritz do? \textcolor{cfared}{[Answer: Issue objective report based on fundamentals]}
			\item What if his boss pressures him to issue favorable recommendation? \textcolor{cfared}{[Answer: Must maintain independence regardless of pressure]}
			\item What could Hilton Brokerage have done to prevent this conflict? \textcolor{cfared}{[Answer: Place Metals \& Mining on restricted list]}
			\item Does Fritz's relationship with company managers matter? \textcolor{cfared}{[Answer: Cannot allow relationships to compromise objectivity]}
		\end{enumerate}
		
		\vspace{0.2cm}
		\textcolor{cfablue}{\textbf{CFA Answer}}: Fritz's analysis must be objective and based solely on company fundamentals. Any pressure from other divisions is inappropriate. Conflict could be eliminated by placing company on restricted list.
		
		\vspace{0.1cm}
		\textcolor{cfared}{\textbf{Common Mistake}}: Thinking business relationships or firm pressure justify compromising research objectivity. They do not.
	\end{frame}
	
	% ============================================================
	% SLIDE 22: HARD QUESTION 16 - LOYALTY TO EMPLOYER LIMITS
	% ============================================================
	\begin{frame}{Hard Question 16: Loyalty to Employer Limits}
		\begin{block}{The Scenario (CFA Handbook, p. 134, Example 11)}
			Meredith Rasmussen discovers a hedge fund's reported performance does not reflect a market decline. Her supervisor and compliance officer tell her to "stay away from the issue." She has been unsuccessful in finding a resolution internally.
		\end{block}
		
		\vspace{0.2cm}
		\textbf{Questions:}
		\begin{enumerate}
			\item Does whistleblowing violate Standard IV(A)? \textcolor{cfared}{[Answer: NO]}
			\item When is whistleblowing justified? \textcolor{cfared}{[Answer: When intent is to protect clients or market integrity]}
			\item What steps must Rasmussen take first? \textcolor{cfared}{[Answer: Follow firm's reporting procedures]}
			\item What if the employer fires her for whistleblowing? \textcolor{cfared}{[Answer: May be protected by law; Standard IV(A) does NOT prohibit justified whistleblowing]}
		\end{enumerate}
		
		\vspace{0.2cm}
		\textcolor{cfablue}{\textbf{CFA Answer}}: Does NOT violate Standard IV(A). Whistleblowing is justified when intent is to protect clients or market integrity, not for personal gain. Must follow firm's reporting procedures first.
		
		\vspace{0.1cm}
		\textcolor{cfared}{\textbf{Common Mistake}}: Thinking whistleblowing always violates loyalty to employer. Protecting market integrity and clients takes precedence.
	\end{frame}
	
	% ============================================================
	% SLIDE 23: HARD QUESTION 17 - SUITABILITY COMPLEXITY
	% ============================================================
	\begin{frame}{Hard Question 17: Suitability and Client Understanding}
		\begin{block}{The Scenario (CFA Handbook, p. 110, Example 5)}
			Max Gubler, CIO of a property/casualty insurance subsidiary, wants to improve diversification and returns. The IPS requires highly liquid investments: large-cap equities and government/corporate bonds with AA minimum and 5-year max maturity. Gubler invests 4\% in a venture capital seed fund with a 3-year lockup and 1/3 per year exit option, leaving equity exposure below the upper limit.
		\end{block}
		
		\vspace{0.2cm}
		\textbf{Questions:}
		\begin{enumerate}
			\item Did Gubler violate Standard III(C)? \textcolor{cfared}{[Answer: YES]}
			\item Did Gubler violate Standard III(A)? \textcolor{cfared}{[Answer: YES]}
			\item What if the IPS didn't explicitly prohibit private equity? \textcolor{cfared}{[Answer: IPS specifies "highly liquid" - lockup violates this]}
			\item What if the returns were guaranteed to be superior? \textcolor{cfared}{[Answer: Still violates - must follow client mandates]}
		\end{enumerate}
		
		\vspace{0.2cm}
		\textcolor{cfablue}{\textbf{CFA Answer}}: Violates Standard III(C) - Suitability and Standard III(A) - Loyalty, Prudence, and Care. Private equity with lockup violates IPS requirement for highly liquid investments. Must understand client's business and circumstances.
		
		\vspace{0.1cm}
		\textcolor{cfared}{\textbf{Common Mistake}}: Thinking returns justify deviations from client mandates. Mandates must be followed regardless of return potential.
	\end{frame}
	
	% ============================================================
	% SLIDE 24: HARD QUESTION 18 - FAIR DEALING IN ALLOCATIONS
	% ============================================================
	\begin{frame}{Hard Question 18: Fair Dealing in Allocations}
		\begin{block}{The Scenario (CFA Handbook, p. 99, Example 2)}
			Morgan Jackson manages a bank's commingled fund and pension accounts. Policy: when a new security is placed on the recommended list, the commingled fund buys first, then pension accounts pro rata. For sales, the commingled fund sells first, then pension accounts. For hot issues, only the commingled fund receives shares if insufficient for all accounts. Discretionary accounts have priority over nondiscretionary accounts.
		\end{block}
		
		\vspace{0.2cm}
		\textbf{Questions:}
		\begin{enumerate}
			\item Does Jackson violate Standard III(B)? \textcolor{cfared}{[Answer: YES]}
			\item If the policy is DISCLOSED, does that make it acceptable? \textcolor{cfared}{[Answer: No - disclosure of unfair policy doesn't make it fair]}
			\item What would be a fair allocation procedure? \textcolor{cfared}{[Answer: Pro rata basis, without prioritization]}
			\item What if clients have different fee structures? \textcolor{cfared}{[Answer: Must still treat all clients fairly]}
		\end{enumerate}
		
		\vspace{0.2cm}
		\textcolor{cfablue}{\textbf{CFA Answer}}: Violates Standard III(B). Policy does not treat all clients fairly. Must execute orders on a systematic basis fair to all clients. Disclosure of unfair policy does not make it fair.
		
		\vspace{0.1cm}
		\textcolor{cfared}{\textbf{Common Mistake}}: Thinking disclosure of unfair policies makes them acceptable. It does NOT.
	\end{frame}
	
	% ============================================================
	% SLIDE 25: HARD QUESTION 19 - COMMUNICATION OF RISK
	% ============================================================
	\begin{frame}{Hard Question 19: Communication of Risk and Limitations}
		\begin{block}{The Scenario (CFA Handbook, pp. 177-178, Example 12)}
			Yuri Yakovlev develops a small-cap quantitative strategy that identifies 10-20 illiquid stocks. Backtests show significant risk-adjusted returns. QSC Capital seeds the strategy with US\$10M. After two years, returns exceed benchmark; Sharpe ratio ~1.0. QSC decides to market to large institutional investors. Yakovlev tells marketing the strategy's capacity is limited - beyond US\$100M, returns may be negatively affected. Marketing manager Wellard says this will "muddy the message" and focuses only on the track record.
		\end{block}
		
		\vspace{0.2cm}
		\textbf{Questions:}
		\begin{enumerate}
			\item Did Yakovlev and Wellard violate Standard V(B)? \textcolor{cfared}{[Answer: YES]}
			\item If the fund is far from the US\$100M threshold, is disclosure still required? \textcolor{cfared}{[Answer: YES]}
			\item What if clients don't understand capacity limitations? \textcolor{cfared}{[Answer: Must educate clients or determine if product is suitable]}
			\item What if the limitation was hypothetical and uncertain? \textcolor{cfared}{[Answer: Still must disclose significant limitations known at the time]}
		\end{enumerate}
		
		\vspace{0.2cm}
		\textcolor{cfablue}{\textbf{CFA Answer}}: Violates Standard V(B). Must disclose significant limitations and risks associated with the investment process. Current investors and prospective investors must be made aware of capacity issues.
		
		\vspace{0.1cm}
		\textcolor{cfared}{\textbf{Common Mistake}}: Thinking limitations don't need disclosure until they become problems. Material limitations must be disclosed.
	\end{frame}
	
	% ============================================================
	% SLIDE 26: HARD QUESTION 20 - CROSS-BORDER COMPLIANCE
	% ============================================================
	\begin{frame}{Hard Question 20: Cross-Border Investment Product Distribution}
		\begin{block}{The Scenario (CFA Handbook, pp. 15-16)}
			Amanda Janney works on a team creating a fixed-income hedge fund to be sold globally through the firm's distribution centers. The firm receives interest from Middle Eastern clients seeking investments complying with Islamic law. The marketing materials do NOT specify whether the fund is suitable for Islamic law compliance. Janney is concerned about the reputation of the fund and firm.
		\end{block}
		
		\vspace{0.2cm}
		\textbf{Questions:}
		\begin{enumerate}
			\item What duty does Janney have regarding Islamic law compliance? \textcolor{cfared}{[Answer: Must be aware of cultural/religious laws and requirements]}
			\item If the fund is NOT Sharia-compliant, what must Janney do? \textcolor{cfared}{[Answer: Disclose this fact in marketing materials]}
			\item What if local law does NOT require disclosure of religious compliance? \textcolor{cfared}{[Answer: Code and Standards may require disclosure to prevent misrepresentation]}
			\item What proactive steps should the firm take? \textcolor{cfared}{[Answer: Acknowledge areas where the fund may not be suitable]}
		\end{enumerate}
		
		\vspace{0.2cm}
		\textcolor{cfablue}{\textbf{CFA Answer}}: As financial markets globalize, members must be aware of differences between cultural/religious laws and governmental regulations. Firms should be proactive in acknowledging areas where products may not be suitable for clients.
		
		\vspace{0.1cm}
		\textcolor{cfared}{\textbf{Common Mistake}}: Thinking only local governmental laws matter. Cultural/religious requirements may affect suitability.
	\end{frame}
	
	% ============================================================
	% SLIDE 27: SUMMARY - WHERE YOUNG IS CORRECT
	% ============================================================
	\begin{frame}{Where Young Is Correct}
		\begin{enumerate}
			\item \textbf{Regulatory Lag}: Laws inevitably lag behind innovation. Financial products evolve faster than regulations.
			\item \textbf{Law as Baseline}: Standard I(A) requires following the "more strict" standard - legal is the floor, not the ceiling.
			\item \textbf{Principled Dissociation}: Sometimes ethical obligations require leaving employment.
			\item \textbf{Civil Disobedience}: History shows laws can be unjust (e.g., apartheid, Jim Crow).
		\end{enumerate}
		\vspace{0.2cm}
		\begin{center}
			\fcolorbox{cfablue}{white}{\parbox{0.8\textwidth}{\centering \textbf{Young is valid as an ANALYTICAL correction to oversimplified reasoning}}}
		\end{center}
	\end{frame}
	
	% ============================================================
	% SLIDE 28: SUMMARY - WHERE YOUNG REQUIRES CAUTION
	% ============================================================
	\begin{frame}{Where Young Requires Caution}
		\begin{enumerate}
			\item \textbf{Risk of Rationalization}: Principled resistance can become self-serving justification.
			\item \textbf{The Misconduct Standard}: Standard I(D) applies regardless of legality. Dishonesty and fraud are always unethical.
			\item \textbf{Public vs. Private Defiance}: Principled defiance is \textcolor{cfablue}{public} and accepts consequences. Quiet noncompliance is not the same.
			\item \textbf{Democratic Legitimacy}: Hughes argues it's "presumptively immoral" for firms to unilaterally decide which laws to respect.
		\end{enumerate}
		\vspace{0.2cm}
		\begin{center}
			\fcolorbox{cfablue}{white}{\parbox{0.8\textwidth}{\centering \textbf{Young's framework must be applied with CARE, not convenience}}}
		\end{center}
	\end{frame}
	
	% ============================================================
	% SLIDE 29: KEY LESSONS FROM HARD QUESTIONS
	% ============================================================
	\begin{frame}{Key Lessons from the 20 Hard Questions}
		\begin{enumerate}
			\item \textbf{Legal is not Ethical}: Compliance is the floor, not the ceiling
			\item \textbf{Disclosure is NOT a Shield}: Disclosing unethical behavior doesn't make it ethical
			\item \textbf{Independence is Paramount}: Cannot be compromised by gifts, pressure, or relationships
			\item \textbf{Client Interests ALWAYS Come First}: Before employer, before personal interests
			\item \textbf{Material Nonpublic Info is Absolute}: No trading, no causing others to trade
			\item \textbf{Due Diligence is Required}: Reasonable basis is non-negotiable
			\item \textbf{Record Retention is Mandatory}: Records are firm property
			\item \textbf{Supervision Requires Action}: Cannot delegate away responsibility
		\end{enumerate}
	\end{frame}
	
	% ============================================================
	% SLIDE 30: FRAMEWORK FOR ETHICAL DECISION-MAKING
	% ============================================================
	\begin{frame}{A Framework for Ethical Decision-Making}
		\small
		\begin{enumerate}
			\item \textbf{Identify Stakeholders} - Who is affected?
			\item \textbf{Evaluate Legal Requirements} - What does the law require?
			\item \textbf{Apply the Code \& Standards} - What do ethical standards require?
			\item \textbf{Consider the "More Strict" Standard} - Which obligation is higher?
			\item \textbf{Assess Impact on Market Integrity} - How does the action affect trust?
			\item \textbf{Document the Analysis} - Could you defend this publicly?
		\end{enumerate}
		
		\vspace{0.2cm}
		\begin{block}{The Three Tests}
			\begin{enumerate}
				\item \textbf{Transparency Test}: Would you defend this action publicly?
				\item \textbf{Social Contract Test}: Does this serve the ultimate benefit of society?
				\item \textbf{Stakeholder Test}: Have you considered impact on all stakeholders?
			\end{enumerate}
		\end{block}
	\end{frame}
	
	% ============================================================
	% SLIDE 31: CONCLUSION
	% ============================================================
	\begin{frame}{Conclusion}
		\begin{block}{Key Takeaway}
			"Although the investment world has become a far more complex place since the first publication of the Standards of Practice Handbook, distinguishing right from wrong remains the paramount principle of the Code and Standards."
			- CFA Institute (2014, p. 5)
		\end{block}
		
		\vspace{0.2cm}
		\textbf{For DBA Scholars:}
		\begin{itemize}
			\item Ethical leadership requires \textcolor{cfablue}{more than legal compliance}
			\item Global business ethics faces challenges laws alone cannot address
			\item Culture of integrity must permeate all levels of operations
			\item Legal standards tell you the \textcolor{cfablue}{floor}, not where to stand
		\end{itemize}
	\end{frame}
	
	% ============================================================
	% SLIDE 32: REFERENCES
	% ============================================================
	\begin{frame}{References}
		\small
		\begin{thebibliography}{9}
			\bibitem{brenkert2019}
			Brenkert, G. G. (2019). Mind the gap! The challenges and limits of (global) business ethics. \textit{Journal of Business Ethics, 155}(4), 917-930.
			
			\bibitem{cfai2014}
			CFA Institute. (2014). \textit{Standards of Practice Handbook} (11th ed.). CFA Institute.
			
			\bibitem{friedman1970}
			Friedman, M. (1970). The social responsibility of business is to increase its profits. \textit{New York Times Magazine}, September 13, 32-33, 122-124.
			
			\bibitem{young2019}
			Young, C. (2019). Putting the law in its place: Business ethics and the assumption that illegal implies unethical. \textit{Journal of Business Ethics, 160}(1), 35-51.
			
			\bibitem{wordpress2013}
			WordPress. (2013, March 15). Ethical theory and its application to contemporary business practice. https://ncys82.wordpress.com/2013/03/15/ethical-theory-and-its-application-to-contemporary-business-practice/
		\end{thebibliography}
	\end{frame}
	
	% ============================================================
	% SLIDE 33: THANK YOU
	% ============================================================
	\begin{frame}{Thank You}
		\begin{center}
			\Huge \textcolor{cfablue}{Questions?}
			
			\vspace{0.5cm}
			\Large \textcolor{cfagold}{"Distinguishing right from wrong remains the paramount principle"}
			
			\vspace{0.5cm}
			\normalsize CFA Institute \textit{Standards of Practice Handbook} (2014, p. 5)
		\end{center}
	\end{frame}
	
	% ============================================================
\end{document}